


Large Language Models (LLMs) are increasingly deployed to autonomously solve real-world tasks. 
A key ingredient for this is the \emph{LLM Function-Calling} paradigm, a widely used approach for equipping LLMs with tool-use capabilities.
However, an LLM calling functions incorrectly can have severe implications, especially when their effects are irreversible, e.g., transferring money or deleting data.
Hence, it is of paramount importance to consider the LLM's confidence that a function call solves the task correctly prior to executing it.
Uncertainty Quantification (UQ) methods can be used to quantify this confidence and prevent potentially incorrect function calls.
In this work, we present what is, to our knowledge, the first evaluation of UQ methods for LLM Function-Calling (FC).
While multi-sample UQ methods, such as Semantic Entropy, show strong performance for natural language Q\&A tasks, we find that in the FC setting the entropy-based multi-sample methods we evaluate offer no clear advantage over simple single-sample UQ methods.
Additionally, we find that the particularities of FC outputs can be leveraged to improve the performance of existing UQ methods in this setting.
Specifically, multi-sample UQ methods benefit from clustering FC outputs based on their abstract syntax tree parsing, while for single-sample UQ methods, restricting the logit-based score to semantically meaningful tokens leaves discrimination essentially unchanged but improves calibration and, by smaller margins, the utility of the score under abstention.

